
\chapter{Нормы регулятора}\label{ch:regulation}
В~платёжном продукте право задаёт архитектуру, когда продукт принимает деньги,
открывает кошелёк, маршрутизирует транзакции или гарантирует продавцу срок
зачисления. Оно определяет роли участников, разделение ответственности,
состояния перевода и~обязательные требования к~безопасности; код либо
реализует эти требования, либо нарушает их.

Глава разбирает российское регулирование; сравнение национальных сетей и~стран
СНГ~--- в~главе~\ref{ch:regional-schemes}. Точные пороги и~даты сверяйте
по~действующим нормативным актам, правилам платёжных систем и~актуальным
материалам участников рынка.

\section{161-ФЗ: субъекты перевода}\label{sec:regulation-161fz}

Федеральный закон от~27.06.2011~№~161-ФЗ \enquote{О~национальной платёжной
системе} (далее~--- 161-ФЗ) отвечает на~вопрос: \emph{кто имеет право менять
записи} в~реестрах банков~\cite{fz-161}. Сеть, процессинг, банк и~платёжный
продукт~--- юридически разные сущности с~разной ответственностью; 161-ФЗ
закрепляет это разделение.

\subsection*{Оператор по переводу денежных средств}

Перевод денежных средств выполняет оператор по~переводу денежных средств
(ОПДС)~--- банк или небанковская кредитная организация (НКО), которые вправе
принять распоряжение к~исполнению, списать деньги у~плательщика и~обеспечить их
зачисление получателю~\cite{fz-161}. У~этих участников~--- эмитента, эквайера,
банка получателя, кошелька с~клиентским остатком (роли в~карточной экосистеме
разбираются в~гл.~\ref{ch:ecosystem} и~гл.~\ref{ch:issuer-acquirer})~---
возникают денежные обязательства и~право изменять денежное состояние.

Пока система только передаёт сообщения и~остаётся интерфейсом к~чужой денежной
функции, статус ОПДС остаётся у~партнёра. Как только продукт ведёт клиентские
остатки, инициирует списание от~своего имени или обещает пользователю хранение
средств, вопрос о~собственном статусе ОПДС нужно решать ещё до~запуска.

\subsection*{Платёжная система и~её оператор}

Закон отдельно выделяет платёжную систему и~оператора платёжной системы.
\highlight{Оператор задаёт правила участия, порядок обмена сообщениями, модель
расчётов, требования к~бесперебойности и~управлению рисками}~\cite{fz-161}.
Жизненный цикл сообщения, обязательные поля, сроки сверки и~правила работы
с~инцидентами~--- следствия правил конкретной системы; в~карточном мире~---
Visa, Mastercard или Мир.

Среди обязательных правил участия 161-ФЗ называет и~порядок приостановления
и~прекращения участия~\cite{fz-161}. После выхода участник сохраняет уже
возникшие обязательства: незавершённые расчёты и~споры по~операциям
до~прекращения закрываются в~сроки, установленные правилами системы. Это
сетевой аналог ответственности эквайера перед уходящим ТСП
(гл.~\ref{ch:issuer-acquirer}).

Провайдер платёжных услуг (payment service provider, PSP), платёжный шлюз или
платёжный оркестратор (маршрутизатор трафика между несколькими PSP
и~эквайерами) упрощают интеграцию, но не~меняют применимый свод правил. Если
сеть требует возможности отмены, фиксированного набора полей или обязательной
связки с~клиринговой записью, продукт исполняет этот свод правил; собственные
соглашения его не~подменяют.

\subsection*{Три функции платёжной инфраструктуры}

\highlight{161-ФЗ разделяет услуги платёжной инфраструктуры на~операционный
  центр (маршрутизирует и~передаёт сообщения между участниками), платёжный
  клиринговый центр (принимает и~сводит распоряжения, рассчитывает нетто-позиции)
  и~расчётный центр (исполняет окончательное движение денег по~счетам
участников)}~\cite{fz-161}. Даже если эти роли совмещены у~одного участника,
маршрутизация авторизации, обработка клиринговых записей и~окончательное
движение денег остаются разными функциями. Сообщение проходит операционный
центр, а~расчёт происходит позже и~через другой механизм.

Система, которая маршрутизирует сообщения, не~обязана отвечать за~расчёт; она
работает на~уровне сообщения. Продукт, обещающий окончательное зачисление
в~определённый момент, опирается на~передачу сообщений и~на~регламент
расчёта~--- со~всеми его сроками и~ограничениями. Национальная система
платёжных карт (НСПК) обеспечивает маршрутизацию и~клиринг внутрироссийского
карточного трафика, но банки остаются эмитентами, эквайерами и~участниками
расчётов.

161-ФЗ закрепляет роли, а~не протоколы. Протокол авторизации меняется,
обязательство банка перед клиентом~--- нет. Закон, который фиксировал бы
технологический стек, требовал бы поправки на~каждый новый стандарт; поэтому
161-ФЗ закрепляет ответственность, а~технологическую реализацию оставляет
положениям Банка России и~правилам платёжных систем, которые обновляются
быстрее закона.

\subsection*{Сопоставление с~PSD2}

В~ЕС эту же задачу разделения ролей решает Directive (EU) 2015/2366 Revised
Payment Services Directive, PSD2~\cite{psd2}. Она делит участников рынка
на~payment service provider (PSP), payment institution, e-money institution,
account servicing payment service provider (ASPSP). Российский ОПДС
соответствует ASPSP плюс части функций payment institution; российский оператор
платёжной системы~--- payment system operator в~терминологии PSD2.

\needspace{4\baselineskip}
PSD2 добавляет два блока требований, которых нет в~161-ФЗ. Во-первых, она
лицензирует роли третьих лиц с~обязательным доступом к~счёту клиента: payment
initiation service provider (PISP) инициирует платёж от~имени
плательщика через счёт в~банке-владельце, account information service
provider (AISP) агрегирует информацию о~счетах клиента. У~161-ФЗ
прямого аналога для~PISP/AISP нет; доступ финтеха к~счёту клиента
в~российской модели строится по~договору с~банком, а~не
по~регуляторному обязательству банка предоставить API. Во-вторых,
PSD2 устанавливает обязательную strong customer authentication (SCA)
на~двух из~трёх факторов (знание, владение, неотъемлемая
характеристика); в~российской регуляторике сопоставимые требования
к~аутентификации распределены по~161-ФЗ (статья~9 об~уведомлении
клиента о~каждой операции) и~положениям Банка России. Подробный
разбор PSD2 и~российской модели открытого банкинга~---
в~главе~\ref{ch:open-banking}.

\section{Положения Банка России}\label{sec:regulation-cbr-provisions}

Положения Банка России переводят требования закона в~правила работы системы:
какие реквизиты обязаны быть в~распоряжении, как защищается платёжная
инфраструктура, какие требования действуют для~банковской информационной
безопасности (ИБ), какие документы получает клиент по карточной операции.

По~состоянию на~март 2026~года инженерные требования по~ИБ задают Положения:
№~821-П~--- защита информации при~переводах, №~850-П~--- операционная
надёжность, №~851-П~--- противодействие переводам без~согласия
клиента~\cite{cbr-821p,cbr-850p,cbr-851p}. Все три отсылают к~ГОСТ~Р~57580.1
в~части состава мер защиты и~к~57580.2 в~части оценки соответствия (разбор
стандарта~--- в~разделе~\ref{sec:regulation-57580}).

\subsection*{762-П: жизненный цикл распоряжения}

Положение Банка России от~29.06.2021~№~762-П \enquote{О~правилах осуществления
перевода денежных средств} (далее~--- 762-П) задаёт реквизиты распоряжения,
порядок приёма и~исполнения, условия отзыва, прохождение через
банки-посредники~\cite{cbr-762p}. У~любого перевода есть машина состояний:
распоряжение получено, принято к~исполнению, исполнено, отклонено, возвращено
или отозвано.

762-П выделяет стадии: приём, удостоверение права распоряжения, авторизация,
исполнение, уведомление~\cite{cbr-762p}. На~примере карточной
покупки за~3\,000~руб. (полный жизненный цикл карточной транзакции
с~авторизацией, клирингом и~расчётом~--- в~гл.~\ref{ch:tx-lifecycle}):

\begin{itemize}
  \item \textbf{T+0, 12:00:00}, приём~--- банк-эмитент получает авторизационный
    запрос от~сети. Момент фиксируется в~журнале; с~этой секунды
    распоряжение считается поступившим.
  \item \textbf{T+0, 12:00:01}, удостоверение права~--- хост эмитента проверяет
    PIN или криптограмму, сверяет статус карты.
  \item \textbf{T+0, 12:00:02}, авторизация~--- проверка достаточности средств,
    скоринг антифрода, решение (одобрение или отказ).
  \item \textbf{T+1\ldots T+2}, клиринг и~исполнение~--- эквайер передаёт
    клиринговую запись, сеть рассчитывает нетто-позиции, расчётный центр
    списывает и~зачисляет.
  \item \textbf{T+2}, уведомление~--- банк-эмитент формирует выписку; продавец
    получает подтверждение зачисления.
\end{itemize}

Для большинства переводов 762-П устанавливает исполнение в~пределах
операционного дня или следующего рабочего дня.

До~принятия распоряжения к~исполнению плательщик вправе его отозвать. После
начала исполнения~--- только пока расчёт не~завершён. Статус \enquote{принято
к~исполнению}~--- юридическая точка, после которой логика \enquote{отмены}
меняется на~логику \enquote{возврата}: другие сроки, другие документы,
обязательное согласование сторон.

Платёжный API, бэк-офис и~внутренняя учётная модель обязаны хранить полную
машину состояний перевода с~временн\'{ы}ми метками каждой стадии. Без места
для~обязательных реквизитов, времени приёма и~времени исполнения система
не~подтверждает правовой статус перевода и~усложняет операционное
сопровождение.

\subsection*{821-П: защита информации при переводе}

Положение Банка России от~17.08.2023~№~821-П \enquote{О~защите информации
при~переводах денежных средств}\footnote{В~силе с~01.04.2024.} задаёт
действующие требования к~защите переводов~\cite{cbr-821p}. Платёжная
инфраструктура становится контролируемой средой с~требованиями к~защищённым
каналам связи, разграничению доступа, журналированию, управлению инцидентами,
оценке соответствия и~средствам криптографической защиты.

Процессинговый коммутатор и~платёжный шлюз нельзя проектировать по~схеме
\enquote{сначала обычный веб-сервис, потом межсетевой экран уровня приложений
(Web Application Firewall, WAF) сверху}. Сегментацию, ключи, администрирование
и~логи планируют на~этапе проектирования.

Сегментация сети постфактум~--- переработка топологии под~рабочей нагрузкой;
замена алгоритмов криптозащиты после запуска~--- миграция данных и~ротация
ключей; разграничение привилегий после выдачи десятков ролей~--- пересмотр
модели доступа. Каждое такое действие дешевле предусмотреть до~запуска, чем
внедрять в~работающей системе.

Надзорная проверка по~821-П опирается на~четыре артефакта, которые Банк России
запрашивает по~графику и~в~ходе проверок.

\begin{itemize}
  \item\textbf{Журнал инцидентов} фиксирует события защиты
    информации с~точностью до~конкретного действия и~принятого
    решения; значимые инциденты передаются в~ФинЦЕРТ по~форме
    0403203\footnote{Утверждена Указанием Банка России
    от~12.01.2022~№~6060-У.}~\cite{cbr-6060u}.
  \item\textbf{Ежегодная самооценка} соответствия требованиям защиты
    информации оформляется по~форме 0409071 Указания Банка России
    от~10.04.2023~№~6406-У \enquote{О~формах отчётности кредитных
    организаций} и~направляется в~Банк России~\cite{cbr-6406u}.
  \item\textbf{Отчётность по~операционной надёжности}~--- сведения
    о~показателях технологических процессов и~применяемых ИТ~---
    сдаётся по~форме 0409072 (того же 6406-У) на~основании Положения
    Банка России от~13.01.2025~№~850-П \enquote{Об~операционной
    надёжности}\footnote{В~силе с~03.05.2025; заменило 787-П.}~\cite{cbr-850p}.
  \item\textbf{Предписания} Банк России выдаёт по~результатам проверок;
    типовые формулировки~--- расхождение декларации и~реализации
    по~средствам криптозащиты, отсутствие журналирования критических
    действий, отсутствующая модель угроз для~конкретного объекта
    инфраструктуры, неполная сегментация сети передачи карточных
    данных. На~каждое предписание организация отвечает планом
    устранения с~привязкой к~конкретному артефакту, и~план
    остаётся под~контролем до~закрытия предписания.
\end{itemize}

\subsection*{851-П: противодействие переводам без согласия клиента}

Положение Банка России от~30.01.2025~№~851-П \enquote{О~противодействии
переводам без согласия клиента}\footnote{В~силе с~29.03.2025; заменило 683-П.}
разделяет требования на~две группы~\cite{cbr-851p}. \emph{Целевое назначение}
первой группы~--- противодействие фроду на~стороне банка-эмитента
(см.~гл.~\ref{ch:fraud}): мониторинг операций по~признакам, установленным
Банком России, информирование клиента, ограничения на~дистанционные операции
по~заявлению клиента, противодействие выдаче кредитных средств под~влиянием
обмана. \emph{Исполнительные меры} охватывают систему ИБ кредитной организации,
требования к~прикладному ПО банковских операций, оценку соответствия
по~57580.2, использование сертифицированных СКЗИ.

\textbf{Прикладное ПО.} Пункт~4.1 851-П распространяется на~всё ПО,
используемое для~осуществления банковских операций: оно должно либо пройти
сертификацию ФСТЭК (по~Постановлению Правительства Российской Федерации
  от~26.06.1995~№~608 \enquote{О~сертификации средств защиты
информации}~\cite{pp-608}), либо иметь оценку соответствия по~оценочному уровню
доверия не~ниже ОУД4 (п.~7.6 раздела~7
ГОСТ~Р~ИСО/МЭК~15408-3-2013~\cite{gost-15408-3}). Так требование
к~прикладному банковскому ПО становится самостоятельным объектом контроля,
прямого аналога которому в~PCI~DSS нет~--- разбор пересечения с~PCI~DSS
см.~в~разделе~\ref{sec:pci-russian}.

\textbf{Уровень соответствия по~57580.2.} Пункт~13 851-П задаёт конкретные
пороги: для~кредитной организации~--- уровень соответствия не~ниже четвёртого
(подп.~\enquote{д} п.~6.9 раздела~6 57580.2); для~филиала иностранного
банка~--- не~ниже третьего (подп.~\enquote{г}) до~31.12.2026 и~не~ниже
четвёртого с~01.01.2027. Внешний независимый аудит проводит проверяющая
организация с~лицензией ФСТЭК, отчёт хранится не~менее пяти лет,
периодичность~--- не~реже одного раза в~два года.

\textbf{Новое относительно 683-П.} 851-П добавляет регистрацию аутентификации
с~геолокацией и~IP-адресами, контроль изменения SIM-карт клиента, уведомление
законных представителей при~выдаче электронных средств платежа (ЭСП)
несовершеннолетним 14--18~лет, модель СУИБ, описанную через цикл PDCA.
С~01.10.2025 для~ряда операций обязательно использование усиленной
неквалифицированной электронной подписи с~подтверждением соответствия ФСБ.

Соответствие PCI~DSS и~соответствие 821-П, 851-П, 850-П~--- параллельные
обязательства, не~альтернативы. Карточная среда, соответствующая PCI~DSS,
остаётся под~российскими регуляторными требованиями: действуют требования 821-П
к~переводам, 851-П к~прикладному ПО и~антифроду, 850-П к~операционной
надёжности. Формулировка \enquote{Мы соответствуем PCI~DSS} подтверждает
выполнение одного свода требований и~не~подтверждает выполнение остальных.

\subsection*{ГОСТ~Р~57580.1 и~57580.2: исполнительный стандарт защиты
информации}\label{sec:regulation-57580}

\highlight{821-П и~851-П не~перечисляют конкретные меры защиты: они~отсылают
  к~национальному стандарту ГОСТ~Р~57580.1-2017, где зафиксирован базовый
состав мер}. Стандарт разработан Банком России совместно с~НПФ
\enquote{КРИСТАЛЛ}~\cite{gost-57580-1}. Парный ГОСТ~Р~57580.2-2018
\enquote{Методика оценки соответствия} задаёт способ внешнего аудита
по~57580.1~\cite{gost-57580-2}.

ГОСТ становится требованием только в~той части, на~которую ссылается
нормативный акт регулятора~--- через нормативную ссылку по~статье~27
Федерального закона от~29.06.2015~№~162-ФЗ \enquote{О~стандартизации
в~Российской Федерации}~\cite{fz-162}. 821-П делает такие ссылки для~конкретных
категорий объектов информатизации; 851-П распространяет требования на~систему
ИБ кредитной организации в~целом, задавая через~57580.2 конкретный минимальный
уровень соответствия (см.~подсекцию выше). За~пределами охвата
обоих Положений 57580 остаётся рекомендательным.

\textbf{Структура 57580.1.} Базовый состав мер разделён на~восемь процессов:
управление доступом, защита вычислительных сетей, контроль целостности
и~защищённости информационной инфраструктуры, защита от~вредоносного кода,
предотвращение утечек, управление инцидентами, защита виртуализации, защита
при~удалённом и~мобильном доступе. К~процессам добавлены четыре направления
управления (планирование, реализация, контроль, совершенствование) и~отдельный
раздел требований к~защите на~этапах жизненного цикла автоматизированных
систем (АС).

\textbf{Три уровня защиты.} Стандарт определяет уровни защиты информации УЗ-1
(усиленный), УЗ-2 (стандартный) и~УЗ-3 (минимальный). Конкретный уровень
для~конкретной АС организация не~выбирает~--- он задаётся нормативным актом
Банка России для~соответствующей категории организации и~типа объекта
информатизации.

\textbf{Оценка соответствия по~57580.2.} Это внешний независимый аудит, а~не
самооценка: его проводит сторонняя организация с~лицензией ФСТЭК России
на~техническую защиту конфиденциальной информации (минимум один из~четырёх
  видов работ: контроль защищённости от~НСД, проектирование защищённых систем,
установка средств защиты, монтаж и~испытания)~\cite{gost-57580-2}. Каждая мера
получает числовую оценку 0 или~1 (выбор), для~процессов управления и~жизненного
цикла АС~--- 0, 0{,}5 или~1 (полнота реализации). Среднеарифметическое
по~процессу даёт итоговый показатель, который сопоставляется с~шестибалльной
качественной шкалой: нулевой уровень (меры не~реализуются), первый
(бессистемно, эпизодически), второй (значительный объём без~контроля), третий
(значительный объём с~эпизодическим контролем), четвёртый (полный объём
со~стабильным контролем), пятый (полный объём с~непрерывным
совершенствованием). Для~кредитной организации 851-П требует уровень
соответствия не~ниже четвёртого; для~филиала иностранного банка~--- не~ниже
третьего до~31.12.2026 и~не~ниже четвёртого с~01.01.2027.

\subsection*{266-П: карточная операция как документ}

Эмиссию платёжных карт и~операции с~их использованием регулирует Положение
Банка России от~24.12.2004~№~266-П \enquote{Об~эмиссии платёжных
карт}~\cite{cbr-266p}; обязанность информировать клиента о~каждой операции
с~использованием ЭСП дополнительно установлена статьёй~9 161-ФЗ~\cite{fz-161}.
На~момент подготовки главы Банк России готовит акт на~смену 266-П: в~2024~году
опубликован проект Положения \enquote{О~порядке предоставления электронных
  средств платежа и~осуществления операций с~их
использованием}.\footnote{По~состоянию на~май 2026~года заменяющий акт
  не~вступил в~силу; реквизиты вступления в~силу проверяйте по~первоисточнику
на~сайте Банка России.} До~его вступления в~силу 266-П сохраняет прямое
действие. Помимо авторизационного жизненного цикла, у~карточного платежа есть
документальный след: чек или иной документ по~операции,
электронное подтверждение, запись в~мобильном банке, уведомление клиенту,
механизм оспаривания.

266-П задаёт состав документа по~операции: наименование документа, номер
и~дату, сумму и~валюту, вид операции, реквизиты карты, наименование
торгово-сервисного предприятия (ТСП) или банкомата, место проведения~--- точный
набор зависит от~канала и~вида операции~\cite{cbr-266p}. Каждая операция должна
однозначно идентифицироваться по~документу, выданному или направленному
клиенту.

Минимальный набор реквизитов на~уровне каждой транзакции~--- исходный материал
для~ответа клиенту, аудитору и~регулятору. Терминальное ПО и~страница оплаты
формируют и~выдают документ с~правильным составом полей. Мобильный банк
и~личный кабинет показывают детали операции в~формате, пригодном
для~оспаривания.

\section{Лицензирование}\label{sec:regulation-licensing}

\subsection*{Интерфейс с~банком-партнёром}

Продукт, который предоставляет пользователю или продавцу интерфейс, но не~берёт
обязательства хранить или переводить деньги, строится в~партнёрстве с~банком,
эквайером или PSP. На~стороне партнёра остаётся денежный статус, расчёты,
регуляторная отчётность и~формальное исполнение распоряжения. На~сам продукт
всё равно распространяются платёжные требования и~требования по~защите
информации через договоры и~интеграции, но самостоятельным носителем банковской
функции он не~становится.

До~запуска архитектура должна ответить на~пять вопросов: где хранятся деньги,
чьё имя стоит в~договоре с~клиентом, кто выпускает чек, кто отвечает за~спорную
операцию, кто меняет баланс. Пока ответы не~зафиксированы, формула \enquote{мы
только интерфейс} не~отвечает ни на~один из~них.

\subsection*{Агрегатор, агент и~PayFac: пограничные модели}

Между чистым интерфейсом и~полноценным хранением денег есть промежуточные
модели:

\begin{itemize}
  \item \textbf{Платёжный агент} действует от~имени банка по~агентскому
    договору по~Федеральному закону от~03.06.2009~№~103-ФЗ
    \enquote{О~деятельности по~приёму платежей физических лиц,
    осуществляемой платёжными агентами}~\cite{fz-103}.
  \item \textbf{Агрегатор} принимает платежи в~пользу нескольких ТСП через
    расчётный счёт банка-партнёра.
  \item \textbf{Платёжный фасилитатор} (payment facilitator, PayFac)~--- модель
    платёжных систем, при которой фасилитатор самостоятельно подключает
    субмерчантов под собственным контрактом с~эквайером.
  \item Иной продукт, который подключает ТСП и~управляет клиентским сценарием,
    хотя денежная функция и~расчётная обязанность остаются у~банка или
    другого партнёра.
\end{itemize}

Модели различаются ответами на~три вопроса: у~кого находятся деньги, кто
выпускает расчётный документ, кто входит в~договорные отношения с~клиентом
и~продавцом.

Подключение ТСП, идентификация клиентов и~ТСП, чек, возврат, спорная операция
и~расчётные сроки распределяются между участниками экосистемы
из~раздела~\ref{sec:ecosystem-participants}: кто заключает договор с~ТСП, кто
несёт риск, кто меняет баланс, кто отвечает за~расчёт после операции. Чем
больше продукт определяет ценообразование, маршрутизацию, поддержку клиентов
и~операции ТСП, тем раньше нужно зафиксировать, где заканчивается сервис
и~начинается регулируемая денежная функция. Слова \enquote{агрегатор},
\enquote{платформа} или \enquote{сервис для~приёма платежей} не~определяют
статус~--- его определяет фактический набор функций.

\subsection*{Процессинг и~маршрутизация без хранения денег}

Платформа, которая маршрутизирует сообщения и~обрабатывает авторизации, но
не~становится должником клиента по~его деньгам, ближе всего к~роли
операционного центра. Если она начинает подтверждать и~сводить распоряжения,
она получает функции платёжного клирингового центра. В~обоих случаях это услуги
платёжной инфраструктуры; деньги остаются у~партнёра~\cite{fz-161}.

Такая система не~превращается автоматически в~банк, но получает требования
к~бесперебойности, защите информации, аудиту и~управлению инцидентами. Формула
\enquote{мы только передаём сообщения} снижает лицензионные требования
по~сравнению с~хранением денег, но~не~освобождает от~инженерной дисциплины.

\subsection*{Кошелёк или баланс клиента}

Продукт, который начинает хранить клиентские средства, вести баланс, открывать
кошелёк или иным образом становиться должником пользователя по~деньгам,
требует статуса кредитной организации (по~Федеральному закону
от~02.12.1990~№~395-1 \enquote{О~банках и~банковской деятельности}) с~правом
на~соответствующие операции, включая перевод электронных денежных средств
(ЭДС)~\cite{fz-161,fz-395-1}. С~этого момента меняются ключевые требования:
капитал, отчётность, комплаенс, антифрод, противодействие отмыванию доходов
и~финансированию терроризма (ПОД/ФТ), безопасность и~отношения с~регулятором.

Денежная функция возникает там, где у~продукта появляется обязательство вернуть
клиенту его деньги по~требованию. Балансы в~личном кабинете, бонусные кошельки
с~возможностью оплаты у~произвольного продавца, предоплаченные счета
маркетплейса с~остатком, конвертируемым в~деньги к~выводу,~--- во~всех этих
случаях продукт хранит остаток клиента и~обязан вернуть его по~запросу. Если
продукт даёт пользователю обещание \enquote{ваши деньги хранятся у~нас
и~доступны по~запросу}, название продукта в~личном кабинете не~меняет правовой
природы операции. Архитектурно остаются два варианта~--- становиться кредитной
организацией самому или встраиваться в~банковскую инфраструктуру через
ОЭДС-партнёра, который остаётся обязанным по~ЭДС (часть~2 статьи~12
  161-ФЗ~\cite{fz-161}; лимиты остатка и~операций по~ЭДС
см.~гл.~\ref{ch:kyc-kyb}). Получение лицензии банка или НКО~--- годы; работа
через ОЭДС-партнёра по~статусу платёжного агрегатора (статья~14.1 161-ФЗ)~---
кварталы, но партнёр определяет тарифную сетку и~учётную модель кошелька,
а~продукт работает в~его операционном дне. Федеральный закон
от~20.02.2026~№~44-ФЗ \enquote{Об~ограничении идентификации физлиц банковскими
платёжными агентами}\footnote{В~силе с~03.03.2026; договоры ОПДС с~БПА
приводятся в~соответствие в~течение 180~дней.}~\cite{fz-44} исключил схему,
в~которой банковский платёжный агент (БПА) самостоятельно идентифицировал
физлицо и~открывал ему электронный кошелёк: эта функция централизована
у~кредитных организаций, и~агентская модель остаётся для~приёма платежей
и~переводов без~открытия счёта в~пределах установленных лимитов.

\subsection*{Собственная платёжная система}

Организация, которая сама задаёт правила взаимодействия участников, обязана
зарегистрироваться оператором платёжной системы по~161-ФЗ~\cite{fz-161}. Вместе
с~регистрацией приходят обязанности по~своду правил, управлению рисками,
непрерывности работы, допуску участников и~модели расчётов.

Пока платформа помогает передавать сообщения, она~--- поставщик технологии. Как
только она задаёт обязательные правила обработки этих сообщений для~нескольких
участников, к~ней применяются обязанности оператора платёжной системы.

\practicenote{%
  Главный архитектурный вопрос~--- возникает~ли у~нас собственное
  обязательство по~чужим деньгам. Если нет, схему можно строить через
  партнёра; если да, схема переходит к~банку, НКО и~полному набору
  регуляторных требований.

}
\section{Пять вопросов к~платёжной архитектуре}\label{sec:regulation-checklist}

Пять вопросов к~новой платёжной схеме или партнёрской интеграции выявляют
правовой, операционный и~продуктовый риск до~запуска. Каждый из~них
в~российской практике 2022--2026 годов уже проявлялся в~публичных случаях.

\subsection*{Где возникает денежное обязательство?}

Юридический должник перед пользователем или продавцом не~всегда совпадает
с~брендом, который видит клиент. Когда Банк России приказом
от~21.02.2024~№~ОД-266 отозвал лицензию у~КИВИ~Банка (финмониторинг
и~разбор причин~--- в~разделе~\ref{sec:aml-qiwi}), держатели вкладов получили выплаты
АСВ до~1{,}4~миллиона рублей, а~держатели QIWI-кошельков~--- нет: остатки
на~кошельках учитывались как электронные денежные средства по~161-ФЗ
и~под~систему страхования вкладов
не~подпадают~\cite{cbr-qiwi-license-2024,fz-161}. Та же схема воспроизводится
в~маркетплейсах: остаток продавца на~внутреннем счёте маркетплейса~---
требование к~маркетплейсу, а~не вклад в~банке-партнёре, и~при банкротстве
маркетплейса страхование вкладов АСВ~не~применяется.

\subsection*{Кто вправе принять распоряжение и~поменять баланс?}

Интерфейс, процессинг и~маршрутизация юридически отделяются от~стороны
с~денежной функцией; перепутать эти роли~--- значит работать без~требуемой
лицензии. С~1~октября 2024~года платёжные агенты обязаны быть включены в~реестр
операторов по~приёму платежей Банка России~\cite{fz-103,cbr-opp-register};
сервисы, маршрутизировавшие платежи между ролью \enquote{агент} и~ролью
\enquote{оператор ЭДС}, включились в~реестр или прекратили работу в~прежней
модели. Тот же критерий определяет, нужна~ли стартапу банковская лицензия: если
продукт сам меняет балансы пользователей, лицензия нужна.

\subsection*{По какому своду правил исполняется операция?}

Закон, правила платёжной системы, договор с~партнёром и~внутренняя модель
состояний не~должны противоречить друг другу. После 6~марта 2022~года карты
Visa и~Mastercard, выпущенные российскими банками, продолжили приниматься
внутри страны через ОПКЦ НСПК; для~эквайеров и~ТСП переход не~потребовал
замены интеграций~--- НСПК с~2015~года маршрутизировала внутрироссийский трафик
по~формату ISO~8583 и~правилам, наследующим Visa~Core~Rules
и~Mastercard~Rules. Администрирование правил для~внутрироссийского трафика
перешло к~НСПК: тарифы межбанковской комиссии (\textit{interchange}) теперь
устанавливает НСПК~\cite{rbc-nspk-fees-2022}, а~интеграция эквайера и~договоры
с~ТСП сохранили прежнюю основу.

\subsection*{Какие статусы, документы и~уведомления система обязана выдавать?}

Если ответа на~этот вопрос нет, сбой проявится на~споре, возврате и~аудите.
СБП~--- это \textit{a2a}-сервис (account-to-account) с~прямым переводом между
счетами клиентов в~разных банках, и~любой межбанковский спор по~такому переводу
с~30~июня 2023~года идёт через систему урегулирования \enquote{Диспут+}
НСПК~\cite{nspk-dispute-plus-sbp-2023} (механика клиентских споров и~отличие
от~карточного цикла~--- в~разделе~\ref{sec:disputes-sbp}). Регламент задаёт
сроки (180~дней на~инициирование, до~30~дней на~рассмотрение), роли инициатора
и~ответчика, состав подтверждающих документов; двусторонние договорённости
между банками его не~подменяют. Если решение положительное, НСПК выступает
арбитром и~проводит автоматические расчёты между сторонами диспута~---
возмещение приходит отдельной проводкой, не~отзывом исходного перевода.

\subsection*{В~какой момент операция становится окончательной и~кто
разбирает спор?}

Безотзывность платежа~--- инженерное свойство, ограниченное правом. Поправки
в~161-ФЗ Федеральным законом от~24.07.2023~№~369-ФЗ \enquote{О~защите клиентов
от~переводов без~согласия}~\cite{fz-369-2023} ввели обязательную сверку
получателя с~базой данных Банка России о~получателях, использованных
в~мошеннических операциях: при~совпадении банк приостанавливает перевод на~срок
до~двух операционных дней до~подтверждения клиентом. Приказ Банка России
от~05.11.2025~№~ОД-2506 \enquote{О~признаках перевода без~согласия
клиента}\footnote{В~силе с~01.01.2026; заменил ОД-1027
от~27.06.2024.}~\cite{cbr-od-2506,fz-161} задаёт критерии, по~которым
банк-эмитент обязан запросить подтверждение клиента и~отказать в~исполнении,
если подтверждение не~получено. Если банк не~приостановил перевод по~получателю
из~базы и~операция оказалась мошеннической, обязан возместить клиенту сумму
в~срок до~30~дней. Для~СБП это означает, что между приёмом распоряжения
и~межбанковской проводкой появляется стадия проверки, в~которой перевод либо
подтверждается клиентом, либо отклоняется банком; финальность по-прежнему
наступает в~момент проводки (детальная механика антифрода и~финальности
в~СБП~--- в~разделе~\ref{sec:sbp-finality}).

На~этапе раннего проектирования без~ответа на~\enquote{кто меняет баланс},
\enquote{кто владеет сводом правил} и~\enquote{кто разбирает ошибочный платёж}
архитектура не~позволяет выполнить обязательства перед регулятором и~клиентом.
На~этапе интеграции, когда API, договор, сценарий поддержки и~внутренняя модель
статусов дают разные ответы на~один и~тот же вопрос, расхождение проявится
на~возврате, диспуте, аудите или инциденте.

\section{Мир: нормативное требование как инженерная
задача}\label{sec:regulation-mir-mandatory}

Закон и~нормативные акты требуют поддержки Мир; для~инженерии это означает
корректное распознавание карт Мир, маршрутизацию внутрироссийского трафика
через инфраструктуру НСПК, сценарии приёма в~терминалах и~электронной
коммерции, отражение Мир в~процедуре подключения ТСП и~в~клиентской
поддержке~\cite{cbr-ps,cbr-nps-faq}.

Инженерная работа PSP или эквайера по~обязательной поддержке Мир состоит
из~терминальной, хостовой и~договорной частей.

На~терминалах~--- обновление конфигурации системы управления терминалами
(Terminal Management System, TMS), развёртывание ядра бесконтактного приёма
Мир, проверка, что терминал распознаёт диапазоны банковских идентификационных
номеров (Bank Identification Number, BIN) карт Мир и~корректно обрабатывает
контактный и~бесконтактный сценарии.

На~хостовой стороне~--- подключение маршрутизации через операционный платёжный
и~клиринговый центр (ОПКЦ), согласование профиля полей формата сообщений
финансовых транзакций ISO~8583 (см.~гл.~\ref{ch:iso8583}) по~спецификации НСПК,
прохождение хостовой сертификации.

В~договорах~--- включение Мир в~стандартный контракт ТСП, настройка
сопоставления кодов категории торгового предприятия (Merchant Category Code,
MCC), корректное отображение в~расчётных отчётах и~на панели управления.

Подробный разбор НСПК, ОПКЦ, Mir Pay и~обязательности Мир~---
в~главе~\ref{ch:mir}. Точные пороги обязательного приёма, даты развёртывания
и~актуальные исключения нужно сверять по~действующим документам НСПК, Банка
России и~правилам сетей.
